%==============================================================================
% sections/02_setup_and_notation.tex
%==============================================================================
\section{Setup, Notation, and Dependence Classes}
\label{sec:setup}

\subsection{Basic objects}

Let \(X=(X_1,\ldots,X_N)\) be an \(N\)-dimensional real-valued loss vector and
\[
  S_N=\sum_{i=1}^N X_i, \qquad M_N=\max_{1\le i\le N} X_i.
\]
Positive values are losses. The target rare-event probability is
\[
  \mu(x)=\Pb(S_N>x), \qquad x\to\infty.
\]
For a vector \(a\in\RR^d\) and factor return vector \(F\in\RR^d\), the
portfolio factor loss is \(L=a^\top F\) after signs have been chosen so that
large positive \(L\) is a bad portfolio outcome. A residual component may be
included as an additional coordinate or block.

A measurable deterministic tie-breaking rule is fixed once and for all. Let
\(R_i(X)=1\) denote the event that coordinate \(i\) is the selected maximum
among \(X_1,\ldots,X_N\). For example, ties can be broken by smallest index.
For each \(i\), write
\[
  S_{-i}=\sum_{j\ne i}X_j, \qquad
  M_{-i}=\max_{j\ne i}X_j,
\]
with \(M_{-i}=-\infty\) if \(N=1\). The conditioning threshold for selected
maximum decompositions is
\[
  T_i(x)=\{x-S_{-i}\}\vee M_{-i},
\]
with strict or non-strict inequalities adjusted according to the deterministic
tie rule. For continuous distributions this distinction is immaterial.

\subsection{Regular variation and common reference tails}

A positive tail \(\Gb\) is regularly varying with index \(-\alpha\), written
\(\Gb\in \RV_{-\alpha}\), if
\[
  \frac{\Gb(tx)}{\Gb(x)}\to t^{-\alpha}, \qquad t>0.
\]
The independent baseline assumes pointwise tail equivalence
\[
  \Pb(X_i>x)\sim c_i\Gb(x), \qquad c_i\ge 0,
\]
for a common \(\Gb\in\RV_{-\alpha}\). The active set is
\[
  \Aset=\{i:c_i>0\}, \qquad C=\sum_{i\in\Aset}c_i.
\]
The financial signed-tail constants are obtained by applying this condition to
the right tail of each signed contribution. If \(X_i=b_iF_i\), then, for a
marginal factor with two-sided tail constants \(p_i^+,p_i^-\),
\[
  \Pb(b_iF_i>x)\sim |b_i|^\alpha
  \{p_i^+\1(b_i>0)+p_i^-\1(b_i<0)\}\Gb(x).
\]

\subsection{Second-order marginal tails}

For second-order expansions we use a direct second-order equivalence. For
active \(i\), there exists \(A(x)\to0\), eventually of constant sign, and
numbers \(d_i\), such that locally uniformly for \(u\) in bounded sets,
\[
  \frac{\Pb(X_i>x-u)}{\Gb(x)}
  = c_i\left(1+\frac{\alpha u}{x}\right)
    + c_id_iA(x) + o(A(x)+x^{-1}).
\]
This formulation is deliberately local because the second-order term used in
the finite-mean expansion shifts the large coordinate by the remaining sum.
Inactive or lighter-tail components are assumed negligible at the retained
second-order scale when the second-order theorem is invoked.

\subsection{Dependence classes}

The expanded paper uses five nested or related dependence classes.

\begin{description}
\item[Independent coordinates.] The \(X_i\)'s are independent and have common
reference regularly varying tails. This is the original model.

\item[Arbitrary dependence with conditional kernels.] The joint law of \(X\) is
arbitrary, but conditional survival kernels
\[
  p_i(t;X_{-i})=\Pb(X_i>t\mid X_{-i})
\]
can be evaluated, approximated, or bounded. This class gives exact
unbiasedness identities but not automatic efficiency.

\item[Latent independent shocks.] Observed factors are generated by
\[
  F=BZ, \qquad Z_1,\ldots,Z_K\ \text{independent},
\]
with regularly varying shock tails. A portfolio loss is \(a^\top F=q^\top Z\),
where \(q=B^\top a\). Dependence among observed factors is inherited from the
mixing matrix \(B\).

\item[Block dependence.] The coordinates are partitioned into blocks
\(B_1,\ldots,B_K\). Blocks are independent; within-block dependence is allowed.
If \(Y_k=\sum_{i\in B_k}X_i\) has a regularly varying tail, the independent
sum theory applies to \(Y_1,\ldots,Y_K\).

\item[Multivariate regular variation.] The full vector satisfies
\[
  \frac{\Pb(X/x\in\cdot)}{\Hb(x)}\xrightarrow{v}\nu(\cdot)
\]
on a cone away from the origin, where \(\nu(tA)=t^{-\alpha}\nu(A)\). In polar
coordinates \(X=R\Theta\), the tail measure decomposes as
\(\alpha r^{-\alpha-1}\dd r\,S(\dd\theta)\). Axis concentration recovers the
independent first-order formula; non-axis spectral mass represents tail
dependence.

\item[Hidden regular variation.] Ordinary MRV may be axis-concentrated while a
slower tail scale exists on pair or block cones. Hidden regular variation is
used to quantify joint-tail regions that are invisible under the first MRV
limit but relevant for second-order risk.
\end{description}

\subsection{Financial loss mapping}

Let \(r_{p,t}\) be portfolio return and \(f_t\) the vector of traded or
constructed factors. A rolling factor model is
\[
  r_{p,t}=\beta_t^\top f_t+\varepsilon_t.
\]
The one-day loss is \(L_t=-r_{p,t}\). Conditional on rolling estimates, the
loss decomposition is
\[
  L_t=\sum_{j=1}^d X_{j,t}+X_{\varepsilon,t},
  \qquad
  X_{j,t}=-\beta_{j,t}f_{j,t}, \quad
  X_{\varepsilon,t}=-\varepsilon_t.
\]
The independent baseline treats the signed factor and residual contributions as
independent tail summands. The dependent extensions either transform the factor
vector into latent shocks, group factors into blocks, estimate a copula
conditional kernel, or estimate the empirical spectral measure of the joint
loss-contribution vector.

\subsection{Boundary and continuity conventions}

For MRV statements, a risk set \(A\) is admissible when it is bounded away from
the origin and \(\nu(\partial A)=0\). For a linear loss \(\ell(z)=a^\top z\),
the set \(A_\ell=\{z:\ell(z)>1\}\) is admissible whenever the spectral measure
assigns zero mass to \(\{\theta:a^\top\theta=1/r\}\) in the relevant polar
normalization, equivalently whenever \(\nu\{z:a^\top z=1\}=0\). In empirical
work we avoid exact-boundary claims by reporting threshold-stability bands.

\subsection{Standing assumptions}
\label{subsec:standing-assumptions}

The following assumptions are invoked throughout the paper. Each main theorem
states which subset it requires.

\begin{assumption}[Common reference regular variation]
\label{ass:rv-common-ref}
There exists a survival function \(\Gb\) with \(\Gb\in\RV_{-\alpha}\) for some
\(\alpha>0\) and constants \(c_i\ge 0\) such that
\(\Pb(X_i>x)\sim c_i\Gb(x)\) as \(x\to\infty\) for every \(i=1,\ldots,N\). When
the first moment is required we assume \(\alpha>1\).
\end{assumption}

\begin{assumption}[Active set non-empty]
\label{ass:active-nonempty}
The active set \(\Aset=\{i:c_i>0\}\) is non-empty and \(C=\sum_{i\in\Aset}c_i>0\).
\end{assumption}

\begin{assumption}[Tie-breaking rule]
\label{ass:tie}
A measurable deterministic tie-breaking rule has been fixed; the events
\(\{R_i(X)=1\}\) partition the sample space. For continuous joint laws the rule
is immaterial.
\end{assumption}

\begin{assumption}[Measurability and integrability]
\label{ass:measurability}
The loss vector \(X\) is a Borel-measurable random vector on a complete
probability space. Conditional survival kernels \(p_i(t;X_{-i})\), where
invoked, are jointly measurable in their arguments.
\end{assumption}

\begin{assumption}[Second-order tail equivalence]
\label{ass:sote}
For each \(i\in\Aset\) there exists an auxiliary function \(A(x)\to 0\),
eventually of constant sign, and a number \(d_i\in\RR\) such that, locally
uniformly in bounded \(u\),
\[
  \frac{\Pb(X_i>x-u)}{\Gb(x)}
  = c_i\left(1+\frac{\alpha u}{x}\right)
    + c_i d_i A(x) + o(A(x)+x^{-1}).
\]
Inactive components are negligible at the retained second-order scale.
\end{assumption}

\Cref{ass:rv-common-ref,ass:active-nonempty,ass:tie,ass:measurability} are the
\emph{independent baseline} assumptions; \cref{ass:sote} is added for
second-order theorems. Dependence-specific assumptions (latent-shock structure,
MRV, hidden regular variation) are introduced in \cref{sec:dependent-cdmcs},
\cref{sec:mrv-spectral}, and \cref{sec:hidden-rv} respectively.
